\documentclass[11pt]{article}
\usepackage[margin=1in]{geometry}
\usepackage{amsmath,amssymb,amsthm}
\usepackage{xurl}
\usepackage{hyperref}
\sloppy
\let\oldus\_ \renewcommand{\_}{\oldus\allowbreak}
\newtheorem{theorem}{Theorem}
\newtheorem{lemma}[theorem]{Lemma}
\newtheorem{corollary}[theorem]{Corollary}
\theoremstyle{definition}
\newtheorem{definition}[theorem]{Definition}
\theoremstyle{remark}
\newtheorem{remark}[theorem]{Remark}

\title{Disproof of conjecture 00000000810:\\ no metric star minimises the first Kirchhoff eigenvalue}
\author{Jackmeson1}
\date{2026-10-05}

\begin{document}
\maketitle

\section{The conjecture and the readings refuted}

\textbf{Statement (English, verbatim).} Definition: Kirchhoff eigenvalues of quantum graphs are the
spectra of metric graphs with Kirchhoff conditions. Conjecture: Among graphs of fixed total length and
vertex count, the minimum of the first eigenvalue is attained by the star graph; and the extremizer is
strictly unique for more than three vertices. (The Chinese text states the same.)

\textbf{Reading.} The conjecture is a conjunction; we refute its first conjunct (``the minimum is
attained by the star graph'') for every vertex count $n\ge 4$ and every total length $L>0$, which
refutes the conjunction. For $n\le 3$ the star $K_{1,n-1}$ is itself a path, so nothing is claimed
there. ``First eigenvalue'' is read as the smallest \emph{positive} Kirchhoff eigenvalue $\lambda_1$
(the eigenvalue $0$, with constant eigenfunctions, is shared by every graph). ``The star graph'' is
read as the metric star $K_{1,n-1}$ with \emph{any} positive edge lengths summing to $L$ (the
equilateral star included). We refute two versions:
\begin{itemize}
\item[(R1)] the competitors are all metric graphs with $n$ vertices and total length $L$;
\item[(R2)] the competitors are only the stars $K_{1,n-1}$ of total length $L$. This covers any
reading that excludes vertices of degree two (the path used in (R1) has such vertices).
\end{itemize}
Our competitor in (R1) is the path $P_n$, a simple connected tree, so (R1) is refuted for every
competitor class that contains it (connected graphs, trees, simple graphs, or graphs with loops and
multiple edges). \emph{Not formalized:} the reading in which ``first eigenvalue'' means the bottom
eigenvalue $0$. Under it the minimum value $0$ is attained by every graph, so uniqueness fails; this
needs the non-negativity of the Kirchhoff Laplacian, which we do not formalize.

\textbf{How the formal statement implies the refutation.} If $\lambda_1(S)\le\lambda_1(G)$ for every
competitor $G$, then $\lambda_1(S)\le\lambda$ for every positive eigenvalue $\lambda$ of every
competitor, because $\lambda_1(G)$ is at most each positive eigenvalue of $G$. The Lean theorems
refute this consequence, so they refute the conjecture. No discreteness of the spectrum is used.

\section{Definitions (as formalized)}

\begin{definition}[metric graph]
A metric graph on the vertex set $\{0,\dots,n-1\}$ (\texttt{MetricGraph n}) consists of finitely
many edges $e$, each with a tail vertex, a head vertex and a length $\ell_e>0$. Edge $e$ is the
interval $[0,\ell_e]$; its endpoint $0$ is glued to the tail and $\ell_e$ to the head. Loops and
multiple edges are allowed. The total length is $\sum_e\ell_e$ (\texttt{totalLength}).
\end{definition}

\begin{definition}[Kirchhoff eigenvalue]
$\lambda\in\mathbb{R}$ is a Kirchhoff eigenvalue (\texttt{IsEigenvalue}) if there are functions $f_e$
with derivatives $f_e'$ such that: (i) $f_e$ has derivative $f_e'(x)$ and $f_e'$ has derivative
$-\lambda f_e(x)$ at every $x\in[0,\ell_e]$, within $[0,\ell_e]$ (one-sided at the endpoints), i.e.\
$-f_e''=\lambda f_e$; (ii) continuity: there are vertex values $F(v)$ with $f_e(0)=F(\mathrm{tail}\,e)$
and $f_e(\ell_e)=F(\mathrm{head}\,e)$; (iii) Kirchhoff: at every vertex $v$,
$\sum_{\mathrm{tail}\,e=v}f_e'(0)-\sum_{\mathrm{head}\,e=v}f_e'(\ell_e)=0$ (the sum of the outgoing
derivatives vanishes; a loop contributes at both ends; at a degree-one vertex this is the Neumann
condition); (iv) $f_e(x)\ne 0$ for some edge $e$ and some $x\in[0,\ell_e]$. \texttt{posEig G} is the
set of positive Kirchhoff eigenvalues.
\end{definition}

\begin{definition}[star and path]
\texttt{star k l}: vertices $0,\dots,k$; edge $e\in\{0,\dots,k-1\}$ joins the centre $0$ to the leaf
$e+1$ and has length $\ell_e$. \texttt{path k L}: vertices $0,\dots,k$; edge $e$ joins $e$ to $e+1$ and
has length $L/k$. Relabelling vertices and edges, or reversing an edge ($x\mapsto\ell_e-x$, which
changes the sign of $f_e'$ and swaps its two ends in (iii)), maps eigenfunctions to eigenfunctions, so
every metric star with $k$ edges has the same Kirchhoff eigenvalues as some \texttt{star k l}.
\end{definition}

\section{Proof}

Throughout, $k=n-1\ge 3$ is the number of edges.

\begin{lemma}[\texttt{ode\_cos}]\label{ode}
Let $s,l>0$. If $f''=-s^2f$ on $[0,l]$ (one-sided derivatives at the ends) and $f'(l)=0$, then
$f(x)=f(l)\cos(s(l-x))$ and $f'(x)=f(l)\,s\sin(s(l-x))$ on $[0,l]$.
\end{lemma}
\begin{proof}
$g=f-f(l)\cos(s(l-\cdot))$ satisfies $g''=-s^2g$ and $g(l)=g'(l)=0$. The energy
$E=g'^2+s^2g^2$ has derivative $2g'(g''+s^2g)=0$, so it is constant on $[0,l]$ (Mathlib
\texttt{constant\_of\_has\_deriv\_right\_zero}), hence $E\equiv E(l)=0$, so $g=g'=0$.
\end{proof}

\begin{lemma}[\texttt{tan\_add\_identity}, \texttt{tan\_sum\_le}, \texttt{tan\_sum\_lt}]\label{tan}
For $x,y\ge 0$ with $x+y<\pi/2$: $\tan x+\tan y+\tan x\tan y\tan(x+y)=\tan(x+y)$. Hence, for
non-negative angles with sum $<\pi/2$, $\sum\tan\theta_i\le\tan\sum\theta_i$; if the angles are
positive and there are at least two of them, the inequality is strict.
\end{lemma}
\begin{proof}
The identity is $\tan(x+y)(1-\tan x\tan y)=\tan x+\tan y$ after clearing the positive cosines. The
inequalities follow by induction on the number of angles, since the product term is $\ge 0$ ($>0$ for
positive angles).
\end{proof}

\begin{lemma}[\texttt{star\_trig}]\label{trig}
Let $k\ge 3$, $\theta_1,\dots,\theta_k>0$ with $\sum\theta_e\le\pi$, and $a_e,c\in\mathbb{R}$ with
$a_e\cos\theta_e=c$ for all $e$ and $\sum_e a_e\sin\theta_e=0$. Then all $a_e=0$.
\end{lemma}
\begin{proof}
Let $\theta_0$ be a largest angle and $R=\sum_{e\ne 0}\theta_e\le\pi-\theta_0$, a sum of at least
two positive terms. For $e\ne 0$, $\theta_e<R\le\pi-\theta_0\le\pi-\theta_e$, so
$\theta_e<\pi/2$, $\cos\theta_e>0$ and $a_e\sin\theta_e=c\tan\theta_e$. Also $\theta_0<\pi$, so
$\sin\theta_0>0$. With $T=\sum_{e\ne 0}\tan\theta_e>0$ the Kirchhoff equation becomes
$a_0\sin\theta_0+cT=0$, and $c=a_0\cos\theta_0$ gives $a_0(\sin\theta_0+T\cos\theta_0)=0$. We show
$P=\sin\theta_0+T\cos\theta_0>0$. If $\cos\theta_0\ge 0$ this is clear. Otherwise
$\theta_0>\pi/2$, and $0<R\le\pi-\theta_0<\pi/2$, so by Lemma~\ref{tan} and the monotonicity of
$\tan$, $T<\tan R\le\tan(\pi-\theta_0)=-\sin\theta_0/\cos\theta_0$. Multiplying by
$\cos\theta_0<0$ gives $T\cos\theta_0>-\sin\theta_0$. So $a_0=0$, then $c=0$, then $a_e=0$ for
$e\ne 0$ since $\cos\theta_e>0$.
\end{proof}

\begin{theorem}[\texttt{star\_no\_eigenvalue\_le}]\label{star}
Let $k\ge 3$ and $\ell_e>0$ with $L=\sum\ell_e$. The star has no Kirchhoff eigenvalue $\lambda$ with
$0<\lambda\le\pi^2/L^2$.
\end{theorem}
\begin{proof}
Let $s=\sqrt\lambda$, so $0<sL\le\pi$. At the leaf $e+1$ only edge $e$ arrives (as its head), so
Kirchhoff gives $f_e'(\ell_e)=0$, and Lemma~\ref{ode} gives $f_e=a_e\cos(s(\ell_e-x))$ with
$a_e=f_e(\ell_e)$. Continuity at the centre gives $a_e\cos(s\ell_e)=F(0)$; Kirchhoff at the centre
gives $\sum_e f_e'(0)=s\sum_e a_e\sin(s\ell_e)=0$. Lemma~\ref{trig} with $\theta_e=s\ell_e$ gives
$a_e=0$ for all $e$, so $f\equiv 0$, contradicting (iv).
\end{proof}

\begin{theorem}[\texttt{path\_eigenvalue}]
The path with $k$ edges of length $L/k$ has total length $L$ and the eigenvalue $\pi^2/L^2$.
\end{theorem}
\begin{proof}
Take $f_e(x)=\cos\big(\tfrac{\pi}{L}(eL/k+x)\big)$, i.e.\ $\cos(\pi t/L)$ in the arclength $t$, with
$F(v)=\cos(\pi v/k)$. The equation holds on each edge. At a vertex $v$ the outgoing derivative of
the edge leaving $v$ is $G(v)=-\tfrac{\pi}{L}\sin(\pi v/k)$, and the derivative of the edge arriving at
$v$ is $G(v)$ too, so the Kirchhoff sum is $G(v)-G(v)=0$ at interior vertices; at the ends only one
term is present and $G(0)=G(k)=0$. Finally $f_0(0)=1\ne 0$.
\end{proof}

\begin{theorem}[\texttt{twoLong\_eigenvalue}, \texttt{star\_beaten}]
For $k=j+2\ge 3$ and lengths $a,a,\varepsilon,\dots,\varepsilon$ ($a,\varepsilon>0$), the star has the
eigenvalue $(\pi/2a)^2$. Consequently, if $\mu>\pi^2/L^2$, some star with $k$ edges and total
length $L$ has a positive eigenvalue $<\mu$.
\end{theorem}
\begin{proof}
With $s=\pi/(2a)$, take $f=\pm\cos(s(a-x))$ on the two long edges and $f=0$ elsewhere. Then
$f=0$ at the centre (since $\cos(sa)=0$), Neumann holds at the leaves, and the outgoing derivatives
at the centre are $s$ and $-s$. For the second claim let $r=\pi/\sqrt\mu<L$, $t=(L+r)/2$, $a=t/2$
and $\varepsilon=(L-t)/j$: then $2a+j\varepsilon=L$ and $(\pi/t)^2<(\pi/r)^2=\mu$.
\end{proof}

\begin{corollary}[\texttt{star\_not\_minimal}, \texttt{not\_star\_minimizer},
\texttt{not\_star\_minimizer\_among\_stars}]
Let $k\ge 3$ and $L>0$. For every star $S$ with $k$ edges and total length $L$, and every
$\mu\in$ \texttt{posEig} $S$: $\pi^2/L^2<\mu$; the path with $k+1$ vertices and length $L$ has
the eigenvalue $\pi^2/L^2$; and another star with $k$ edges and length $L$ has a positive eigenvalue
$<\mu$. Hence there is no star $S$ (any lengths summing to $L$) whose least positive eigenvalue
$\mu$ (\texttt{IsLeast (posEig S) mu}) satisfies $\mu\le\lambda$ for all positive eigenvalues $\lambda$
of all metric graphs with $k+1$ vertices and length $L$ (R1), nor for all positive eigenvalues of all
such stars (R2).
\end{corollary}

\begin{remark}
For context: Friedlander (Ann.\ Inst.\ Fourier 55 (2005) 199--211,
\url{https://doi.org/10.5802/aif.2095}) states in the abstract: ``We establish the sharp lower bound
for eigenvalues of a metric graph.'' Our proof does not use that paper; all steps above are proved in
Lean.
\end{remark}

\section{Lean formalization and axiom audit}

File \texttt{lean/Conjecture810/Basic.lean} (430 lines, namespace \texttt{C810}, only
\texttt{import Mathlib}; Lean 4.33.1, Mathlib v4.33.1). Main theorems:
\begin{verbatim}
theorem not_star_minimizer (k : N) (hk : 3 <= k) (L : R) (hL : 0 < L) :
    not exists (l : Fin k -> R) (hl : forall e, 0 < l e), sum e, l e = L /\
      exists mu, IsLeast (posEig (star k l hl)) mu /\
        forall G : MetricGraph (k + 1), G.totalLength = L ->
          forall lam in posEig G, mu <= lam

theorem not_star_minimizer_among_stars (k : N) (hk : 3 <= k) (L : R)
    (hL : 0 < L) :
    not exists (l : Fin k -> R) (hl : forall e, 0 < l e), sum e, l e = L /\
      exists mu, IsLeast (posEig (star k l hl)) mu /\
        forall (l' : Fin k -> R) (hl' : forall e, 0 < l' e),
          sum e, l' e = L ->
          forall nu in posEig (star k l' hl'), mu <= nu
\end{verbatim}
Supporting theorems: \texttt{star\_no\_eigenvalue\_le}, \texttt{path\_eigenvalue},
\texttt{twoLong\_eigenvalue}, \texttt{star\_beaten}, \texttt{star\_not\_minimal}, \texttt{star\_trig},
\texttt{ode\_cos}, \texttt{tan\_sum\_lt}.
\texttt{\#print axioms} for \texttt{not\_star\_minimizer}, \texttt{not\_star\_minimizer\_among\_stars},
\texttt{star\_not\_minimal}, \texttt{star\_no\_eigenvalue\_le} and \texttt{path\_eigenvalue}:
\texttt{[propext, Classical.choice, Quot.sound]}. There is no \texttt{sorry}, no
\texttt{native\_decide} and no added axiom.

\end{document}
